\chapterimage{figs/chapterhead.png}
\chapter{Standalone Tools}
\label{chap:standalone-tools}

This chapter collects the current standalone tools that sit beside the main
GUI and the shell. Here, \textit{standalone tools} means auxiliary front ends
and support entry points, not the core kernel classes themselves.

\noindent\textbf{Objective.} Present the standalone tools that accompany GenESyS and explain how they fit the user workflow.

\noindent\textbf{Audience.} Users who need auxiliary tools for analysis, optimization, or specialized workflows.

\noindent\textbf{Scope.} Keep the chapter limited to user-facing tools, their intended purpose, and the level of support currently documented. The chapter describes the current build surface conservatively and does not claim full end-to-end certification for every helper window.

\noindent\textbf{Status.} Partial documentation, grounded in the current CMake targets and source tree.

\section{Current standalone front ends}

The current repository exposes four standalone Qt front ends that share the
same backend libraries as the main application family. Figure~\ref{fig:application-executable-matrix}
in Chapter~\ref{chap:application-overview} already maps them at a high level;
this chapter explains what each one is for.

\begin{itemize}
\item \textbf{HTTP Worker GUI.}

\noindent\textbf{Target.} \texttt{genesys\_httpworker\_gui\_application}.

\noindent\textbf{Executable.} \texttt{genesys-httpworker-gui}.

\noindent\textbf{Purpose.} A control dialog for the worker runtime. It is intended for worker-related configuration and interaction, not for general model editing.

\item \textbf{Data Analyser.}

\noindent\textbf{Target.} \texttt{genesys\_dataanalyser\_gui\_application}.

\noindent\textbf{Executable.} \texttt{genesys-dataanalyser-gui}.

\noindent\textbf{Purpose.} A dedicated analysis window for inspecting simulation result data without opening the full model editor.

\item \textbf{Optimizer.}

\noindent\textbf{Target.} \texttt{genesys\_optimizer\_gui\_application}.

\noindent\textbf{Executable.} \texttt{genesys-optimizer-gui}.

\noindent\textbf{Purpose.} A specialized window for optimization-oriented workflows.

\item \textbf{AI Assistant.}

\noindent\textbf{Target.} \texttt{genesys\_ai\_assistant\_gui\_application}.

\noindent\textbf{Executable.} \texttt{genesys-ai-assistant-gui}.

\noindent\textbf{Purpose.} A graphical assistant front end for AI-assisted workflows.
\end{itemize}

These helpers are built as separate targets in the GUI umbrella, and each one
pulls only the backend libraries it needs. In practice, that separation keeps
the main editor lighter and makes the auxiliary workflows easier to reason
about.

\section{Independent GUI processes and Tools menu}
\label{sec:standalone-tools-process}

The main editor's \textbf{Tools} menu contains \textbf{Data Analyzer},
\textbf{Optimization}, and \textbf{AI Assistant}. Each action requests
its own installed executable through \texttt{ToolLauncher} and Qt's
\texttt{QProcess::startDetached}; it does not construct a tool window in
the main GUI process. Each standalone program constructs a separate
\texttt{QApplication}, \texttt{Simulator}, and top-level window. Closing
one of these applications does not terminate the editor.

The \texttt{gui-app} CMake preset builds only the main editor, while
\texttt{gui-dataanalyser}, \texttt{gui-optimizer}, and
\texttt{gui-ai-assistant} remain independent build presets. A development
build of the editor alone therefore does not guarantee that all Tools
menu items are available. The Debian package build enables these three
optional targets and installs them in the installation prefix's
\texttt{bin} directory. The main GUI itself is installed below
\path{libexec/genesys/system/bin}; the launcher service resolves the
documented prefix-relative directory without executing arbitrary
programs from \texttt{PATH}. The portable distribution instead places
the actual tool executables beside the main GUI, with public wrapper
commands under the portable installation's \texttt{bin} directory.

A tool's local \texttt{Simulator} does not automatically contain the
model currently open in the editor. No model pointer crosses the process
boundary, and the Tools actions do not currently implement model
serialization or interprocess state synchronization. Independent
process startup must not be confused with scientific or AI-provider
workflow validation.

The Data Analyzer and Optimizer window implementation files still
reside under \path{source/applications/gui/genesys/tools/} for historical
reasons. Their respective standalone CMake targets compile those
files directly, and the main GUI source target excludes their
translation units. A future focused migration may move the files
into the standalone application directories, but deleting them now
would remove active window implementations. The source inventory and
related validation plan are maintained in
\path{docs/ai_assistants/reference/STANDALONE_TOOLS_GUI.md}.

\section{Support-package entry point}

The repository also contains the \texttt{source/tools} package, which is a
support layer for analysis and numerical workflows. Its small command-line
entry point prints the available tool interfaces with \texttt{--list} or
\texttt{--help}. The output currently covers:

\begin{itemize}
\item analysis interfaces such as \texttt{DataSet\_if} and \texttt{DataAnalyser\_if};
\item fitting and hypothesis-testing interfaces;
\item distribution abstractions;
\item numerical support interfaces, including solver, quadrature, root-finding,
and ODE contracts;
\item optimization and factorial-design helpers.
\end{itemize}

That entry point is mainly a diagnostic inventory for developers and power
users. It confirms what the package exports, but it is not the primary route
for everyday model authoring or simulation work.

\section{Usage boundary}

When the task is interactive model editing or execution, the main GUI remains
the first choice. When the task is scripted execution, use the shell. When the
task is worker control, analysis, optimization, or AI-assisted helper flows,
choose the corresponding standalone tool instead of forcing the main GUI to do
everything.

The practical rule is simple: select the smallest tool that matches the task.
That keeps the workflow clearer for users and keeps the documentation aligned
with the current build layout.
